% !TEX root = hw10-answers-graded.tex
\chead{\textbf{Exercises week 10}}
\externaldocument{../hw9/hw9}
\begin{document}
\flushleft\textbf{Topic: Structural Causal Models: equivalences and marginalisations}

\begin{enumerate}[label=10.\arabic*]
    \item For a simple SCM $M$ and a subset $L \subset V$ of endogenous variables, we have that $G(M_{\setminus L})$ is a subgraph $G(M)^{\setminus L}$. Notably, these graphs are not necessarily equal. Give a simple SCM $M$ and marginalised SCM $M_{\setminus L}$ for some subset of nodes $L\subset V$ for which $G(M_{\setminus L}) \neq G(M)^{\setminus L}$. (Explicitly give the graphs $G(M_{\setminus L})$ and $G(M)^{\setminus L}$.)

    \emph{Hint: you can use exercise \ref*{ex:unfaithful_scm} for inspiration.}
    \begin{ungradedsolution}
        Consider the SCM $M$ with endogenous variables $V = \{X, Y, Z\}$, exogenous variables $W = \{\varepsilon_X, \varepsilon_Y, \varepsilon_Z\}$, structural equation
        \begin{align*}
            X & = \varepsilon_X         \\
            Y & = X + \varepsilon_Y     \\
            Z & = X - Y + \varepsilon_Z
        \end{align*}
        and exogenous distribution $\varepsilon_X, \varepsilon_Y, \varepsilon_Z \sim \Ncal(0, \sigma^2)$ independent for some $\sigma\in\mathbb{R}$. Its graph $G(M)$ and the subsequent marginalisation $G(M)_{\setminus \{Y\}}$ w.r.t.\ $Y$ are given by:
        \begin{figure}[!htb]
            \centering
            \begin{subfigure}[]{0.45\textwidth}
                \centering
                \begin{tikzpicture}[scale=1, transform shape]
                    \node[ndlat] (EX) at (-1, 1.5) {$\varepsilon_X$};
                    \node[ndout] (X) at (0, 0) {$X$};
                    \node[ndlat] (EY) at (1, 3) {$\varepsilon_Y$};
                    \node[ndout] (Y) at (1, 1.5) {$Y$};
                    \node[ndlat] (EZ) at (3, 1.5) {$\varepsilon_Z$};
                    \node[ndout] (Z) at (2, 0) {$Z$};
                    \draw[arlout] (EX) to (X);
                    \draw[arlout] (EY) to (Y);
                    \draw[arlout] (EZ) to (Z);
                    \draw[arout] (X) to (Y);
                    \draw[arout] (Y) to (Z);
                    \draw[arout] (X) to (Z);
                \end{tikzpicture}
                \caption{$G(M)$}
            \end{subfigure}
            \begin{subfigure}[]{0.45\textwidth}
                \centering
                \begin{tikzpicture}[scale=1, transform shape]
                    \node[ndlat] (EX) at (-1, 1.5) {$\varepsilon_X$};
                    \node[ndout] (X) at (0, 0) {$X$};
                    \node[ndlat] (EY) at (1, 3) {$\varepsilon_Y$};
                    % \node[ndout] (Y) at (1, 1.5) {$Y$};
                    \node[ndlat] (EZ) at (3, 1.5) {$\varepsilon_Z$};
                    \node[ndout] (Z) at (2, 0) {$Z$};
                    \draw[arlout] (EX) to (X);
                    \draw[arlout] (EY) to (Z);
                    \draw[arlout] (EZ) to (Z);
                    % \draw[arout] (X) to (Y);
                    % \draw[arout] (Y) to (Z);
                    \draw[arout] (X) to (Z);
                \end{tikzpicture}
                \caption{$G(M)_{\setminus \{Y\}}$}
            \end{subfigure}
            % \caption{Causal graph $G$.}
            % \label{fig:one}
        \end{figure}

        $M$ is simple so it is essentially uniquely solvable wrt.\ $Y$, with $g^{[Y]}(X, \varepsilon_Y) = X + \varepsilon_Y$ the corresponding partial solution function; substituting this into the structural equation for $(X, Z)$ gives the marginalised SCM $M_{\setminus \{Y\}}$ with endogenous variables $V = \{X, Z\}$, exogenous variables $W = \{\varepsilon_X, \varepsilon_Y, \varepsilon_Z\}$, structural equation
        \begin{align*}
            X & = \varepsilon_X                                                           \\
            Z & = X - (X + \varepsilon_Y) + \varepsilon_Z = \varepsilon_Z - \varepsilon_Y
        \end{align*}
        and exogenous distribution $\varepsilon_X, \varepsilon_Y, \varepsilon_Z \sim \Ncal(0, \sigma^2)$ independent. The graph of this marginalised SCM is as follows:
        \begin{figure}[!htb]
            \centering
            \begin{tikzpicture}[scale=1, transform shape]
                \node[ndlat] (EX) at (-1, 1.5) {$\varepsilon_X$};
                \node[ndout] (X) at (0, 0) {$X$};
                \node[ndlat] (EY) at (1, 3) {$\varepsilon_Y$};
                % \node[ndout] (Y) at (1, 1.5) {$Y$};
                \node[ndlat] (EZ) at (3, 1.5) {$\varepsilon_Z$};
                \node[ndout] (Z) at (2, 0) {$Z$};
                \draw[arlout] (EX) to (X);
                \draw[arlout] (EY) to (Z);
                \draw[arlout] (EZ) to (Z);
                % \draw[arout] (X) to (Y);
                % \draw[arout] (Y) to (Z);
                % \draw[arout] (X) to (Z);
            \end{tikzpicture}
            \caption{$G(M_{\setminus \{Y\}})$}
        \end{figure}
    \end{ungradedsolution}

    \item % ($1 + 3\times 2$ points) 
    Prove Proposition \ref*{prp:equivalent_scms}: Let $M = (J, V, W, \mathcal{X}, P, f)$ and $\tilde{M} = (J, V, W, \mathcal{X}, P, \tilde{f})$ be two SCMs which can only differ in their causal mechanisms, but such that $M \equiv \tilde{M}$. Prove the following claims:
    \begin{enumerate}
        \item Equivalence is preserved by interventions:
        \begin{enumerate}
            \item $M_{\doit(T\dots)} \equiv \tilde{M}_{\doit(T\dots)}$ for $T \subseteq V \cup J$ (hard interventions of all three types);
            \item $M_{\doit(T\ot\hat{f}_T)} \equiv \tilde{M}_{\doit(T\ot\hat{f}_T)}$ for $T \subseteq V$ and $\hat f_T : \CX_J \times \CX_W \times \CX_V \to \CX_T$ a measurable function
            (soft interventions);
            \item $M_{\doit(I_B)} \equiv \tilde{M}_{\doit(I_B)}$ for $B \subseteq V \cup J$ (adding intervention variables);
        \end{enumerate}
        % \item Equivalence is preserved by twinning.
        \item Equivalent iSCMs have equivalent submodels and composition preserves equivalence;
        \item Equivalent iSCMs have the same null sets:
        if $M \equiv \tilde{M}$ then $N$ is an $M$-null set if and only if it is an $\tilde{M}$-null set, and `$M\as$' means the same as `$\tilde{M}\as$';
        \item Equivalent iSCMs have the same
        \begin{enumerate}
            \item potential outcomes
            \item solutions
            \item Markov kernels
            \item (partial) solution functions;
        \end{enumerate}
        \item (Partial) solvability, essentially unique (partial) solvability, essential uniqueness of (partial) solution functions and simplicity are invariants under equivalence:
        if $M \equiv \tilde{M}$ then
        \begin{enumerate}
            \item $M$ is solvable w.r.t.\ $L \ins V$ if and only if $\tilde{M}$ is solvable w.r.t.\ $L$,
            \item $M$ is essentially uniquely solvable w.r.t.\ $L \ins V$ if and only if it $\tilde{M}$ is essentially uniquely solvable w.r.t.\ $L$,
            \item $g^{[L]}$ is an essentially unique solution function of $M$ w.r.t.\ $L \ins V$ if and only if it is an essentially unique solution function of $\tilde{M}$ w.r.t.\ $L$,
            \item $M$ is simple if and only if $\tilde{M}$ is simple;
        \end{enumerate}
        \item Essential linearity is preserved.
    \end{enumerate}
    % \begin{enumerate}
    %     \item Prove that equivalence of SCMs ($M \equiv \tilde{M}$) is indeed an equivalence relation.
    %     \item \label{ex:equiv} Equivalence is preserved by hard interventions, i.e.: if $M$ is equivalent to $\tilde{M}$ and $T \subseteq J \cup V \cup W$ is an intervention target, then $M_{\doit(T\dots)}$ is equivalent to $\tilde{M}_{\doit(T\dots)}$. Do this for the three variants of hard interventions.
    %     \item Show that equivalent SCMs have the same solution functions (Definition \ref*{def:solution-function}). Deduce from this and from Exercise \ref{ex:equiv} that unique solvability and simplicity are invariants of equivalence, i.e.: if $M$ is equivalent to $\tilde{M}$ and $M$ is uniquely solvable or simple, then $\tilde{M}$ is also uniquely solvable or simple, respectively.
    %     \item Prove that equivalent SCMs have the same potential outcomes (Definition \ref*{def:scm_solutions_for_inputs}), solutions (Definition \ref*{def:scm_solutions}), and the same Markov kernels (Definition \ref*{def:scm_kernels}).
    % \end{enumerate}
    \begin{ungradedsolution}
        In all parts we assume that $M = \langle J, V, W, \mathcal{X}, P, f \rangle$ and $\tilde{M} = \langle J, V, W, \mathcal{X}, P, \tilde{f} \rangle$ are two SCMs which can only differ in their causal mechanisms.

        \begin{enumerate}
            \item The first three claims follow directly from the definitions.
            \item The first three claims follow directly from the definitions.
            \item The first three claims follow directly from the definitions.
            % Let $M \equiv \tilde{M}$.
            % The condition
            % $$\forall v \in V: x_v = f_v(\xJ,\xV,\xW) \iff x_v = \tilde{f}_v(\xJ,\xV,\xW) \qquad M\as$$
            % implies $S = \tilde{S}\ M\as$, where
            % $$S = \{(x_J,x_W,x_V) \in \CX_{J\cup W\cup V} \mid x_V = f(x_J,x_W,x_V) \},$$
            % $$\tilde{S} = \{(x_J,x_W,x_V) \in \CX_{J\cup W\cup V} \mid x_V = \tilde{f}(x_J,x_W,x_V) \}.$$
            % It even implies that $S \,\Delta\, \tilde{S}$ is measurable.
            % % Indeed, for all $v \in V$, define the sets $S_v := \{x \in \CX : x_v = f_v(x)\}$ and
            % % $\tilde{S}_v := \{x \in \CX : x_v = \tilde{f}_v(x)\}$, the functions
            % % $$h_v : \CXJ \times \CXW \times \CXV \to \CX_v \times \CX_v : (\xJ,\xW,\xV) \mapsto (x_v, f_v(\xV,\xJ,\xW))$$
            % % $$\tilde h_v : \CXJ \times \CXW \times \CXV \to \CX_v \times \CX_v : (\xJ,\xW,\xV) \mapsto (x_v, \tilde f_v(\xV,\xJ,\xW))$$
            % % and the diagonals $\Delta_v = \{(x_v,x_v) : x_v \in \CX_v\} \subseteq \CX_v^2$.
            % % Since $h_v$ is measurable and $\Delta_v$ is a measurable set (because $\CX_v$ is Hausdorff), $S_v = h_v^{-1}(\Delta_v)$ is a measurable set.
            % % Similarly, $\tilde S_v = \tilde h_v^{-1}(\Delta_v)$ is a measurable set.
            % % Hence $S_v \,\Delta\, \tilde S_v$ is measurable.
            % The assumption $M \equiv \tilde M$ then implies that $S_v \,\Delta\, \tilde S_v$ is contained in an $M$-null set of the form $N_v = \tilde{N}_v \times \CX_V$ with $\tilde{N}_v \ins \CX_J \times \CX_W$ a measurable $M$-null set.

            \item We prove the fourth point.
            $M \equiv \tilde{M}$ implies that
            \begin{equation}\label{eq:prp:equivalent_scms}
                x_V = f(x) \iff x_V = \tilde{f}(x)\ M\as.
            \end{equation}

            \begin{enumerate}
                \item
                Potential outcomes are invariant.
                Let $x_J \in \CX_J$.
                We show that if a random variable $x^{\doit(x_J)}_{V,W}$ is a potential outcome of $M$, then it is a potential outcome of $\tilde{M}$.
                Using Lemma~\ref{lem:M-as_to_rv-as}, \eref{eq:prp:equivalent_scms} implies:
                $$X_V^{\doit(\xJ)} = f(x_J,X_W^{\doit(\xJ)},X_V^{\doit(\xJ)}) \iff X_V^{\doit(\xJ)} = \tilde{f}(x_J,X_W^{\doit(\xJ)},X_V^{\doit(\xJ)}) \text{ a.s.}.$$
                Since $X_W^{\doit(\xJ)} \sim P_M(X_W$), this proves the claim.

                \item
                Solutions are invariant.
                Let $(\C{U} \times \CXJ, K(U|X_J))$ be a transition probability space and $X : \C{U} \times \CXJ \to \CX$ be a conditional random variable.
                We show that if $X$ is a solution of $M$, then it is a solution of $\tilde{M}$.
                Using Lemma~\ref{lem:M-as_to_crv-as}, \eref{eq:prp:equivalent_scms} implies:
                $$X_V = f(X_J,X_W,X_V) \iff X_V = \tilde{f}(X_J,X_W,X_V) \text{ a.s.}.$$
                Since $K(X_W|X_J) = P_M(X_W$), this proves the claim.

                \item
                Markov kernels are invariant: this follows from the invariance of solutions.

                \item
                Partial solution functions are invariant.
                Let $L \ins V$.
                We show that if $g^{[L]}$ is a partial solution function of $M$ w.r.t.\ $L$, then it is a partial solution function of $\tilde{M}$ w.r.t\ $L$.
                By assumption, $g^{[L]} : \CXJ \times \CX_{V\sm L} \times \CXW \to \CX_L$ is a measurable function.
                Because $M\as$ and $\tilde{M}\as$ mean the same,
                $$g^{[L]}(\xJ,x_{V\sm L},\xW) = f_L\big(\xJ, x_{V\sm L}, g^{[L]}(\xJ,x_{V\sm L},\xW), \xW\big) \quad M\as$$
                together with \eref{eq:prp:equivalent_scms} implies:
                $$g^{[L]}(\xJ,x_{V\sm L},\xW) = \tilde{f}_L\big(\xJ, x_{V\sm L}, g^{[L]}(\xJ,x_{V\sm L},\xW), \xW\big) \quad \tilde{M}\as,$$
                which proves the claim.
            \end{enumerate}

            \item The fifth point follows straightforwardly because $M$ and $\tilde{M}$ have the same partial solution functions if they are equivalent.

            \item The preservation of essential linearity follows because the coefficients of an affine mapping are identified if the mapping is known up to an $M$-null set.
        \end{enumerate}
        % In all parts we assume that $M = \langle J, V, W, \mathcal{X}, P, f \rangle$ and $\tilde{M} = \langle J, V, W, \mathcal{X}, P, \tilde{f} \rangle$ are two SCMs which can only differ in their causal mechanisms.
        % \begin{enumerate}
        %     \item This follows directly from Definition \ref*{def:scm_obs_interv_equiv} and the fact that logical equivalence ``$\Longleftrightarrow$'' is an equivalence relation.
        %     \item We discuss only the case of hard interventions with specified target value, since the other two cases work the same way (and are even slightly easier to show).
        %     Thus, let $T \subseteq J \cup V \cup W$ be an intervention target and $\xi_{T} \in \mathcal{X}_T$ be the intervention target value.
        %     The intervened SCMs are then given by
        %     \begin{align*}
        %         M_{\doit(X_T = \xi_T)} = \left\langle J \setminus T, V \cup T, W \setminus T, \mathcal{X}, P_{W \setminus T}, (f_{V \setminus T}, \xi_T) \right\rangle, \\
        %         \tilde{M}_{\doit(X_T = \xi_T)} = \left\langle J \setminus T, V \cup T, W \setminus T, \mathcal{X}, P_{W \setminus T}, (\tilde{f}_{V \setminus T}, \xi_T) \right\rangle.
        %     \end{align*}
        %     Now, set $f' := (f_{V \setminus T}, \xi_T)$ and $\tilde{f}' := (\tilde{f}_{V \setminus T}, \xi_T)$.

        %     Then, for $v \in V \setminus T$ and arbitrary $x \in \mathcal{X}$ we obtain:
        %     \begin{align*}
        %         x_v = f'_v(x) & \Longleftrightarrow x_v = f_v(x)           \\
        %                       & \Longleftrightarrow x_v = \tilde{f}_v(x)   \\
        %                       & \Longleftrightarrow x_v = \tilde{f}'_v(x).
        %     \end{align*}
        %     Thereby, in the first step we used the definition of $f'$ and that $v \in V \setminus T$.
        %     In the second step, we used the equivalence $M \equiv \tilde{M}$.
        %     In the last step, we used the definition of $\tilde{f}'$ and again that $v \in V \setminus T$.

        %     Finally, for $v \in T$ and arbitrary $x \in X$, we obtain:
        %     \begin{align*}
        %         x_v = f'_v(x) & \Longleftrightarrow x_v = \xi_v            \\
        %                       & \Longleftrightarrow x_v = \tilde{f}'_v(x).
        %     \end{align*}
        %     Since $V \cup T = (V \setminus T) \cup T$, this proves the claim.
        %     \item Let $M \equiv \tilde{M}$, and let $g: \mathcal{X}_J \times \mathcal{X}_W \to \mathcal{X}_V$ be any measurable function.
        %     Then the following equivalences hold:
        %     \begin{align*}
        %         g \text{ is a solution function for } M & \Longleftrightarrow \forall x_J, \forall x_W: \ \ g(x_J, x_W) = f(x_J, g(x_J, x_W), x_W)                          \\
        %                                                 & \Longleftrightarrow \forall v, \forall x_J, \forall x_W: \ \ g_v(x_J, x_W) = f_v(x_J, g(x_J, x_W), x_W)           \\
        %                                                 & \Longleftrightarrow   \forall v, \forall x_J, \forall x_W: \ \ g_v(x_J, x_W) = \tilde{f}_v(x_J, g(x_J, x_W), x_W) \\
        %                                                 & \Longleftrightarrow   \forall x_J, \forall x_W: \ \ g(x_J, x_W) = \tilde{f}(x_J, g(x_J, x_W), x_W)                \\
        %                                                 & \Longleftrightarrow g \text{ is a solution function for } \tilde{M}.
        %     \end{align*}
        %     This shows that equivalent SCMs have the same solution functions, and consequently that unique solvability is an invariant of equivalence.

        %     Now, assume that $M$ is simple and that $M \equiv \tilde{M}$.
        %     We want to show that $\tilde{M}$ is also simple.
        %     For this, let $T \subseteq J \cup V \cup W$ be any subset.
        %     Then we need to show that $\tilde{M}_{\doit(T)}$ is uniquely solvable.
        %     But note that due to the simplicity of $M$, $M_{\doit(T)}$ is uniquely solvable.
        %     Furthermore, as was shown in Exercise \ref{ex:equiv}, $M_{\doit(T)} \equiv \tilde{M}_{\doit(T)}$. We have just shown that unique solvability is an invariant of equivalence, and thus the unique solvability of $\tilde{M}_{\doit(T)}$ automatically follows.

        %     \item Let $M \equiv \tilde{M}$.
        %     Let $X_{V\cup W}^{\doit(x_J)}$ be any random variable with codomain $\mathcal{X}_V \times \mathcal{X}_W$, and with $x_J \in \mathcal{X}_J$.
        %     We have the following equivalences:
        %     \begin{align*}
        %                             & X_{V \cup W}^{\doit(x_J)} \text{ is a potential outcome of } M \text{ for input } x_J \in \mathcal{X}_J                                         \\
        %         \Longleftrightarrow & X_W^{\doit(x_J)} \sim P \text{ and } X_V^{\doit(x_J)} = f(x_J, X_V^{\doit(x_J)}, X_W^{\doit(x_J)}) \text{ a.s. }                                \\
        %         \Longleftrightarrow & X_W^{\doit(x_J)} \sim P \text{ and } \forall v \in V: \ \ X_v^{\doit(x_J)} = f_v(x_J, X_V^{\doit(x_J)}, X_W^{\doit(x_J)}) \text{ a.s. }         \\
        %         \Longleftrightarrow & X_W^{\doit(x_J)} \sim P \text{ and } \forall v \in V: \ \ X_v^{\doit(x_J)} = \tilde{f}_v(x_J, X_V^{\doit(x_J)}, X_W^{\doit(x_J)}) \text{ a.s. } \\
        %         \Longleftrightarrow & X_W^{\doit(x_J)} \sim P \text{ and } X_V^{\doit(x_J)} = \tilde{f}(x_J, X_V^{\doit(x_J)}, X_W^{\doit(x_J)}) \text{ a.s. }                        \\
        %         \Longleftrightarrow & X_{V \cup W}^{\doit(x_J)} \text{ is a solution for } \tilde{M} \text{ with input } x_J \in \mathcal{X}_J.
        %     \end{align*}
        %     This shows that equivalent SCMs have the same potential outcomes. Since this holds for all $x_J \in \Xcal_J$ and potential outcomes are, by definition, just evaluations of a solution, we have that equivalent SCMs have the same solutions as well.

        %     From Definition \ref*{def:scm_kernels} it is immediate that SCMs with the same solutions have the same Markov kernels as well.
        % \end{enumerate}
    \end{ungradedsolution}

    \item % ($6 \times 1$ points) 
    Given SCM: $M = \left\langle J=\emptyset, V=\{A, B, C, D\}, W=\{A', B', C', D'\}, \R^8, P, f \right\rangle$, where $P$ are four independent normal distributions with means and standard deviations $\mu_A, \sigma_A, \mu_B, ..., \sigma_D$, and
    \[
        \begin{pmatrix}
            f_A(x, \varepsilon) \\ f_B(x, \varepsilon) \\ f_C(x, \varepsilon) \\ f_D(x, \varepsilon)
        \end{pmatrix}
        =
        \begin{pmatrix}
            0           & 0           & 0 & 0           \\
            \alpha_{ab} & 0           & 0 & \alpha_{db} \\
            0           & \alpha_{bc} & 0 & \alpha_{dc} \\
            \alpha_{ad} & 0           & 0 & 0           \\
        \end{pmatrix}
        \begin{pmatrix}
            x_A \\ x_B \\ x_C \\ x_D
        \end{pmatrix}
        +
        \begin{pmatrix}
            \varepsilon_{A'} \\ \varepsilon_{B'} \\ \varepsilon_{C'} \\ \varepsilon_{D'}
        \end{pmatrix}
    \]
    for nonzero real numbers $\alpha_{(\cdot)}$. Note that we have a slightly different notation as in the lecture notes, as we refer to $x_V$ as $x$, and refer to $x_W$ as $\varepsilon$.
    \begin{enumerate}
        \item Construct a solution function.
        \item Give the graph $G^+(M)$.
        % TODO: G^+(M), and also in other exercises
        \item For the set of latent variables $L=\{D\}$, give the marginal SCM $M_{\setminus L}$ and its solution function.
        \item Show that $G^+(M_{\setminus L})$ is a subgraph of $G^+(M)^{\setminus L}$.
        \item Show observational equivalence of $M$ and $M_{\setminus L}$ with respect to $V \setminus L$ directly using the solution functions and Markov kernels.
        \item Show observational equivalence between $M_{\doit(B)}$ and ${\left(M_{\setminus L}\right)}_{\doit(B)}$ with respect to $\{A,C\}$.
    \end{enumerate}

    \begin{gradedsolution}
        \begin{enumerate}
            \item \points{1} As the SCM is acyclic with order e.g. $(A,D,B,C)$, the solution can be read off from $f$. Alternatively, use  $(1_4-B_{VV})^{-1}\varepsilon$ (where $B_{VV}$ relates to the notation in Proposition \ref*{prp:linear_scms}).
            One finds:
            \[\begin{cases}
                    g_A(\varepsilon) = \varepsilon_{A'}                                                                                                                                                                                                                             \\
                    g_B(\varepsilon) = \varepsilon_{A'} \left(\alpha _{ab}+\alpha _{ad} \alpha _{db}\right)+\varepsilon_{B'}+\varepsilon_{D'} \alpha _{db}                                                                                                                          \\
                    g_C(\varepsilon) = \varepsilon_{A'} \left(\alpha _{ab} \alpha _{bc}+\alpha _{ad} \alpha _{bc} \alpha _{db}+\alpha _{ad} \alpha _{dc}\right)+\varepsilon_{B'} \alpha _{bc}+\varepsilon_{D'} \left(\alpha _{bc} \alpha _{db}+\alpha _{dc}\right)+\varepsilon_{C'} \\
                    g_D(\varepsilon) = \varepsilon_{A'} \alpha _{ad}+\varepsilon_{D'}
                \end{cases}\]
            \item \points{0.5} $G(M):$
            \begin{center}
                \begin{tikzpicture}[scale=1, transform shape]
                    \tikzstyle{every node} = [draw,shape=circle];
                    \node (A) at (0, 2) {$A$};
                    \node (B) at (1, 0) {$B$};
                    \node (C) at (3, 0) {$C$};
                    \node (D) at (2, 2) {$D$};
                    \node (epsA) at (0, 4) {${A'}$};
                    \node (epsB) at (1, -2) {${B'}$};
                    \node (epsC) at (3, -2) {${C'}$};
                    \node (epsD) at (2, 4) {${D'}$};
                    \draw[-{Latex[length=3mm, width=2mm]}] (A) to (B);
                    \draw[-{Latex[length=3mm, width=2mm]}] (A) to (D);
                    \draw[-{Latex[length=3mm, width=2mm]}] (D) to (C);
                    \draw[-{Latex[length=3mm, width=2mm]}] (D) to (B);
                    \draw[-{Latex[length=3mm, width=2mm]}] (B) to (C);
                    \draw[-{Latex[length=3mm, width=2mm]}] (epsA) to (A);
                    \draw[-{Latex[length=3mm, width=2mm]}] (epsB) to (B);
                    \draw[-{Latex[length=3mm, width=2mm]}] (epsC) to (C);
                    \draw[-{Latex[length=3mm, width=2mm]}] (epsD) to (D);
                \end{tikzpicture}
            \end{center}

            \item \points{1} Intervening on $\{A,B,C\}$, we solve for $x_D$ and find $g'_D(x_{(A, B, C)}, \varepsilon)=\alpha_{ad} x_A + \varepsilon_{D'}$ as a solution function of $M_{\doit(A, B, C)}$. Filling this in at $f$ leads to:
            \begin{align*}
                \begin{pmatrix}
                    \tilde{f}_A(x) \\ \tilde{f}_B(x) \\ \tilde{f}_C(x)
                \end{pmatrix}
                 & =
                \begin{pmatrix}
                    0           & 0           & 0 & 0           \\
                    \alpha_{ab} & 0           & 0 & \alpha_{db} \\
                    0           & \alpha_{bc} & 0 & \alpha_{dc} \\
                \end{pmatrix}
                \begin{pmatrix}
                    x_A \\ x_B \\ x_C \\ \alpha_{ad} x_A + \varepsilon_{D'}
                \end{pmatrix}
                +
                \begin{pmatrix}
                    \varepsilon_{A'} \\ \varepsilon_{B'} \\ \varepsilon_{C'}
                \end{pmatrix} \\
                 & =
                \begin{pmatrix}
                    0                                  & 0           & 0 \\
                    \alpha_{ab}+\alpha_{ad}\alpha_{db} & 0           & 0 \\
                    \alpha_{ad}\alpha_{dc}             & \alpha_{bc} & 0 \\
                \end{pmatrix}
                \begin{pmatrix}
                    x_A \\ x_B \\ x_C
                \end{pmatrix}
                +
                \begin{pmatrix}
                    \varepsilon_{A'} \\ \varepsilon_{B'} +\alpha_{db}\varepsilon_{D'}\\ \varepsilon_{C'} + \alpha_{dc}\varepsilon_{D'}
                \end{pmatrix}
            \end{align*}
            as the causal mechanism of $M_{\setminus L}$. The solution function of $M_{\setminus L}$ is
            \[\begin{cases}
                    \tilde{g}_A(\varepsilon) = \varepsilon_{A'}                                                                                                    \\
                    \tilde{g}_B(\varepsilon) = \varepsilon_{A'} \left(\alpha _{ab}+\alpha _{ad} \alpha _{db}\right)+\varepsilon_{B'}+\varepsilon_{D'} \alpha _{db} \\
                    \tilde{g}_C(\varepsilon) = \varepsilon_{A'} \left(\alpha _{ab} \alpha _{bc}+\alpha _{ad} \alpha _{bc} \alpha _{db}+\alpha _{ad} \alpha _{dc}\right)+\varepsilon_{B'} \alpha _{bc}+\varepsilon_{D'} \left(\alpha _{bc} \alpha _{db}+\alpha _{dc}\right)+\varepsilon_{C'},
                \end{cases}\]
            which follows directly from Lemma \ref*{lem:marginalization_preserves_unique_solvability}.

            \item \points{0.5} $G(M_{\setminus L}):$
            \begin{center}
                \begin{tikzpicture}[scale=1, transform shape]
                    \tikzstyle{every node} = [draw,shape=circle];
                    \node (A) at (0, 2) {$A$};
                    \node (B) at (1, 0) {$B$};
                    \node (C) at (3, 0) {$C$};
                    \node (epsA) at (0, 4) {${A'}$};
                    \node (epsB) at (1, -2) {${B'}$};
                    \node (epsC) at (3, -2) {${C'}$};
                    \node (epsD) at (2, 4) {${D'}$};
                    \draw[-{Latex[length=3mm, width=2mm]}] (A) to (B);
                    \draw[-{Latex[length=3mm, width=2mm]}] (A) to (B);
                    \draw[-{Latex[length=3mm, width=2mm]}] (A) to (C);
                    \draw[-{Latex[length=3mm, width=2mm]}] (epsD) to (B);
                    \draw[-{Latex[length=3mm, width=2mm]}] (epsD) to (C);
                    \draw[-{Latex[length=3mm, width=2mm]}] (B) to (C);
                    \draw[-{Latex[length=3mm, width=2mm]}] (epsA) to (A);
                    \draw[-{Latex[length=3mm, width=2mm]}] (epsB) to (B);
                    \draw[-{Latex[length=3mm, width=2mm]}] (epsC) to (C);
                \end{tikzpicture}
            \end{center}

            $G(M)^{\setminus L}:$
            \begin{center}
                \begin{tikzpicture}[scale=1, transform shape]
                    \tikzstyle{every node} = [draw,shape=circle];
                    \node (A) at (0, 2) {$A$};
                    \node (B) at (1, 0) {$B$};
                    \node (C) at (3, 0) {$C$};
                    \node (epsA) at (0, 4) {${A'}$};
                    \node (epsB) at (1, -2) {${B'}$};
                    \node (epsC) at (3, -2) {${C'}$};
                    \node (epsD) at (2, 4) {${D'}$};
                    \draw[-{Latex[length=3mm, width=2mm]}] (A) to (B);
                    \draw[-{Latex[length=3mm, width=2mm]}] (A) to (B);
                    \draw[-{Latex[length=3mm, width=2mm]}] (A) to (C);
                    \draw[-{Latex[length=3mm, width=2mm]}] (epsD) to (B);
                    \draw[-{Latex[length=3mm, width=2mm]}] (epsD) to (C);
                    \draw[-{Latex[length=3mm, width=2mm]}] (B) to (C);
                    \draw[-{Latex[length=3mm, width=2mm]}] (epsA) to (A);
                    \draw[-{Latex[length=3mm, width=2mm]}] (epsB) to (B);
                    \draw[-{Latex[length=3mm, width=2mm]}] (epsC) to (C);
                    \draw[{Latex[length=3mm, width=2mm]}-{Latex[length=3mm, width=2mm]}, bend left] (B) to (C);
                \end{tikzpicture}
            \end{center}

            We see that $G(M_{\setminus L})$ is a subgraph of $G(M)^{\setminus L}$.
            \item \points{0.5} \label{ans:obs_equiv} The projections of both solution functions $g$ and $\tilde{g}$ onto $V\setminus L = \{A, B, C\}$ are the same, so the result follows directly from the definition of observable equivalence (Definition \ref*{def:scm_obs_interv_equiv}) and Remark \ref*{rem:using_solution_functions} (if the solutions functions are the same, then the Markov kernels coincide).
            \item \points{0.5} For both SCMs we have the solution function
            \[\begin{cases}
                    g_A(x_B, \varepsilon) = \varepsilon_{A'} \\
                    g_C(x_B, \varepsilon) = \varepsilon_{A'} \alpha_{ad} \alpha _{dc}+x_B \alpha _{bc}+\varepsilon_{C'}+\varepsilon_{D'} \alpha_{dc},
                \end{cases}\]
            so similar to \ref{ans:obs_equiv} we conclude that they are observably equivalent.
        \end{enumerate}
    \end{gradedsolution}

    % \item BONUS: show that $M-as \implies M_{\doit(Y)}-as $ and $M^{[L]}-as$.
    \item \textbf{BONUS:} Construct an iSCM $M$ with $L\subseteq O\subseteq V$ such that the marginalized iSCM $M_{\setminus L}$ is essentially uniquely solvable w.r.t. $O \setminus L$, but the iSCM $M$ is not essentially uniquely solvable w.r.t. $O$.
\end{enumerate}

\end{document}
